\section{The Kabbalah Model in Summary and Comparison}

\subsection{The Whole Model in One Arc}

Here is the whole Kabbalah model in a single arc. Before anything there is the
\textbf{Infinite}, the unknowable source beyond all name. It \textbf{contracts},
withdrawing from a point to leave room for a world, then sends a ray of light back
in. The light organizes into the \textbf{ten sefirot} and descends through
\textbf{four worlds}, the building blocks being the twenty-two letters and the
divine names. The early vessels cannot hold the light and \textbf{shatter},
scattering holy \textbf{sparks} into a broken world of shells. After the break
comes \textbf{repair}, the sefirot reorganizing into stable faces. Into this broken
world is planted the \textbf{human soul}, made in the image of the divine form,
whose deeds can find and free the fallen sparks. The course of history is that
gathering. When the last spark returns, the exile of the \textbf{Shekhinah} ends
and the light shines whole again. Creation is the break. History is the mending.
The human being is the mender.

\subsection{Against the Universal Theme}

This collection runs one thread through every tradition. The \textbf{one goes out into
the many and returns}. Kabbalah is the clearest statement of that thread in the
whole collection. It does not merely gesture at it. It maps every stage.

The going-out is exact. The Infinite withdraws so the many can be, then pours
itself into vessels. The scattering is exact. The vessels break and the one light
falls into countless fragments hidden in countless shells. And the return is
exact. Through the labor of its own creatures, the one gathers its scattered light
back into wholeness. Kabbalah is the universal theme drawn out to full machinery,
with a named source, a named scattering, and a named return.

\begin{longtable}{L{0.36\linewidth} L{0.42\linewidth}}
\caption{The universal theme in plain terms and in Kabbalah terms.}\\
\toprule
\textbf{Plain-term concept} & \textbf{Kabbalah term} \\
\midrule
The one, the source & Ein Sof, the Infinite \\
The going-out & Tzimtzum, the contraction \\
The structure that pours out & The ten sefirot, the four worlds \\
The scattering into the many & Shevirat ha-kelim, the breaking of the vessels \\
The fragments of the one & The holy sparks (nitzotzot) \\
The broken many & The qelippot, the shells \\
The worker who gathers & The human soul \\
The return, the repair & Tikkun, the gathering of the sparks \\
The reunion of the one & The redemption of the Shekhinah \\
\bottomrule
\end{longtable}

\subsection{How It Will Rhyme with the Buddhist Model}

The next model in this collection is the Buddhist one. It rhymes with Kabbalah on
the same deep theme while differing sharply on the mechanism. The rhyme is worth
naming in advance, so the contrast lands when we reach it.

\begin{itemize}
 \item \textbf{Emanation versus dependent arising}. Kabbalah has the many flow
 \textbf{out of} a single source by emanation, a real descent from the one.
 Buddhism has the many \textbf{arise together} by dependent origination,
 with no single source behind them. One model pours out from above. The
 other co-arises with no above.
 \item \textbf{The sparks versus buddha-nature}. Kabbalah locates the hidden
 divine in \textbf{scattered sparks} trapped in the world, fragments to be
 gathered up. Buddhism locates the awakened reality in
 \textbf{buddha-nature}, already whole and present in each being, to be
 uncovered rather than collected.
 \item \textbf{Tikkun versus liberation}. Kabbalah aims at \textbf{repair}, the
 mending of a real cosmic break and the return of the light. Buddhism aims
 at \textbf{liberation}, the release from a turning that was never solid to
 begin with. One fixes a broken cosmos. The other wakes from a dream of
 one.
 \item \textbf{Evil as broken structure versus evil as ignorance}. Kabbalah
 reads evil as \textbf{broken structure}, light fallen into the wrong place,
 with no power of its own. Buddhism reads the root trouble as
 \textbf{ignorance}, a misperception that binds, dissolved by seeing
 clearly.
\end{itemize}

The two models agree that the surface world is a fall from or a veiling of a deeper
wholeness, and that the human task is to undo the separation. They disagree on
whether there was ever a real one to fall, and on whether the cure is to gather
the fragments or to wake from the splitting.

\subsection{How Kabbalah Frames the Goal}

Kabbalah frames the goal as \textbf{repair and return}. The cosmos broke at its
root, the one light scattered into the many shells, and the divine presence fell
into exile with it. The goal is not escape from the world and not the erasure of
the self. It is the \textbf{mending of the break} from inside the world. The human
soul, through deeds done with intention, lifts the trapped sparks one by one and
carries them home. When the last spark is gathered, the broken vessels are whole
again, the Shekhinah is reunited with her source, and the scattered light shines
as one. The goal of Kabbalah is exactly this, the repair and return that reunites
the scattered light with its source.
